\documentclass{beamer}
\usepackage[utf8]{inputenc}
\usepackage[portuguese]{babel}
\usepackage{amsmath}
\usepackage{graphicx}
\usepackage{booktabs}
\usepackage{algorithm}
\usepackage{algorithmic}

\usetheme{Madrid}
\usecolortheme{default}

\title[Otimização de Layout]{Otimização de Layout e Maior Retângulo Vazio}
\subtitle{Comparação entre Solução Analítica (MIP) e Metaheurística (Evolução Diferencial)}
\author{Gustavo Pessanha}
\date{\today}

\begin{document}

% --- Slide de Título ---
\begin{frame}
    \titlepage
\end{frame}

% --- Sumário ---
\begin{frame}{Sumário}
    \tableofcontents
\end{frame}

% ============================================================================
% SEÇÃO 1: INTRODUÇÃO
% ============================================================================
\section{Introdução do Problema}

\begin{frame}{O Problema do Maior Retângulo Vazio}
    \textbf{Objetivo:}
    Encontrar a melhor disposição de um conjunto de obstáculos (polígonos e círculos) dentro de uma área de trabalho (canvas) de forma a maximizar a área do maior retângulo livre (vazio) possível.

    \vspace{0.5cm}
    \textbf{Aplicações Práticas:}
    \begin{itemize}
        \item Corte de materiais (guilhotina).
        \item Design de layouts industriais.
        \item Posicionamento de janelas em interfaces gráficas.
    \end{itemize}

    \vspace{0.5cm}
    \textbf{Desafio:}
    É um problema de otimização combinatória e geométrica complexo (NP-difícil), envolvendo empacotamento (Packing) e geometria computacional.
\end{frame}

\begin{frame}{Visualização do Problema}
    \begin{columns}
        \column{0.5\textwidth}
        \textbf{Entrada:}
        \begin{itemize}
            \item Dimensões do Canvas ($W \times H$).
            \item Lista de formas (Polígonos, Círculos).
        \end{itemize}
        
        \vspace{0.5cm}
        \textbf{Saída:}
        \begin{itemize}
            \item Coordenadas $(x_i, y_i)$ para cada forma.
            \item Dimensões $(w, h)$ e posição do maior retângulo vazio.
        \end{itemize}

        \column{0.5\textwidth}
        \centering
        \includegraphics[width=\linewidth]{A_1764733018_optimized.png}
        % Espaço reservado para imagem se houvesse
    \end{columns}
\end{frame}

\begin{frame}{Exemplo Visual de Soluções}
    \centering
    \begin{columns}
        \column{0.33\textwidth}
        \centering
        \includegraphics[width=\linewidth]{A_1764732618_initial.png}
        \par\vspace{0.2cm}
        \tiny \textbf{Estado Inicial}
        
        \column{0.33\textwidth}
        \centering
        \includegraphics[width=\linewidth]{A_1764732618_optimized.png}
        \par\vspace{0.2cm}
        \tiny \textbf{Analítico (MIP)} \\ Área: 708.42
        
        \column{0.33\textwidth}
        \centering
        \includegraphics[width=\linewidth]{M_1764732618_optimized.png}
        \par\vspace{0.2cm}
        \tiny \textbf{Metaheurística (DE)} \\ Área: 739.00
    \end{columns}

    \vspace{0.5cm}
    \footnotesize
    Comparação visual para uma instância com 5 polígonos.
    \begin{itemize}
        \item \textbf{Analítico:} Solução compacta, otimizada para perímetro (proxy).
        \item \textbf{Metaheurística:} Solução alternativa, neste caso encontrando uma área ligeiramente maior.
    \end{itemize}
\end{frame}

% ============================================================================
% SEÇÃO 2: MODELOS MATEMÁTICOS
% ============================================================================
\section{Modelos Matemáticos}

\begin{frame}{Definição Formal}
    Seja $\mathcal{S} = \{S_1, S_2, \dots, S_N\}$ o conjunto de formas a serem posicionadas no canvas $C = [0, W] \times [0, H]$.
    
    \vspace{0.3cm}
    \textbf{Variáveis de Decisão:}
    \begin{itemize}
        \item $\mathbf{p}_i = (x_i, y_i)$: Posição da forma $S_i$.
        \item $R = (x_R, y_R, w, h)$: O retângulo vazio.
    \end{itemize}

    \vspace{0.3cm}
    \textbf{Restrições Gerais:}
    \begin{enumerate}
        \item \textbf{Contenção:} $S_i(\mathbf{p}_i) \subset C, \forall i$.
        \item \textbf{Não-Sobreposição (Formas):} $S_i \cap S_j = \emptyset, \forall i \neq j$.
        \item \textbf{Não-Sobreposição (Retângulo):} $R \cap S_i = \emptyset, \forall i$.
    \end{enumerate}

    \vspace{0.3cm}
    \textbf{Função Objetivo:}
    \[ \text{Maximizar } A(R) = w \cdot h \]
\end{frame}

\begin{frame}{Modelo 1: Programação Inteira Mista (MIP)}
    \textbf{Simplificação:} Todas as formas são aproximadas por suas \textit{Bounding Boxes} (Caixas Envolventes).
    
    \vspace{0.3cm}
    \textbf{Variáveis:}
    \begin{itemize}
        \item $x_i, y_i$: Coordenadas do canto inferior esquerdo da forma $i$.
        \item $b_{ij}^k \in \{0, 1\}$: Variáveis binárias para separação relativa.
    \end{itemize}

    \vspace{0.3cm}
    \textbf{Restrições de Não-Sobreposição (Big-M):}
    Para cada par $(i, j)$, pelo menos uma deve ser ativa:
    \begin{align*}
        x_i + w_i &\le x_j + M(1 - b_{ij}^1) \quad (\text{i à esquerda de j}) \\
        x_j + w_j &\le x_i + M(1 - b_{ij}^2) \quad (\text{j à esquerda de i}) \\
        y_i + h_i &\le y_j + M(1 - b_{ij}^3) \quad (\text{i abaixo de j}) \\
        y_j + h_j &\le y_i + M(1 - b_{ij}^4) \quad (\text{j abaixo de i}) \\
        \sum_{k=1}^4 b_{ij}^k &\ge 1
    \end{align*}
\end{frame}

\begin{frame}{Modelo 1: Função Objetivo Linearizada}
    O solver MIP (CBC) requer funções lineares. A área $w \cdot h$ é quadrática.
    
    \vspace{0.5cm}
    \textbf{Estratégia Adotada:}
    Maximizar o \textbf{Perímetro} como proxy para a Área.
    
    \[ \text{Maximizar } Z = w + h \]
    
    \vspace{0.3cm}
    \textit{Nota:} Embora correlacionado, maximizar o perímetro não garante a maximização da área (ex: retângulo $1 \times 100$ tem perímetro 202 e área 100, enquanto $50 \times 50$ tem perímetro 200 e área 2500).
\end{frame}

\begin{frame}{Modelo 2: Otimização Contínua (Metaheurística)}
    \textbf{Abordagem:} Evolução Diferencial (Differential Evolution).
    
    \vspace{0.3cm}
    \textbf{Codificação:} Vetor real $\mathbf{v} = [x_1, y_1, x_2, y_2, \dots, x_N, y_N]$.
    
    \vspace{0.3cm}
    \textbf{Função de Custo (com Penalidade):}
    \[ F(\mathbf{v}) = \begin{cases} 
        - \text{Area}(R) & \text{se válido} \\
        1000 + 10 \cdot (\text{Area}_{\text{overlap}} + \text{Area}_{\text{out}}) & \text{se inválido}
    \end{cases}
    \]
    
    \vspace{0.3cm}
    \textbf{Cálculo da Área (Inner Problem):}
    \begin{itemize}
        \item Rasterização do canvas em uma grade binária.
        \item Algoritmo de "Maior Retângulo em Histograma" ($O(N)$ por linha).
    \end{itemize}
\end{frame}

% ============================================================================
% SEÇÃO 3: SOLUÇÕES E IMPLEMENTAÇÃO
% ============================================================================
\section{Soluções e Implementação}

\begin{frame}{Implementação Analítica (Python + PuLP)}
    \begin{itemize}
        \item \textbf{Biblioteca:} PuLP com solver CBC.
        \item \textbf{Vantagem:} Garante a otimalidade global para o modelo simplificado (Bounding Box + Perímetro).
        \item \textbf{Desvantagem:} Complexidade exponencial. O número de variáveis binárias cresce quadraticamente com o número de formas ($O(N^2)$).
        \item \textbf{Tempo Limite:} Definido em 60 segundos.
    \end{itemize}
\end{frame}

\begin{frame}{Implementação Metaheurística (Scipy + Shapely)}
    \begin{itemize}
        \item \textbf{Algoritmo:} \texttt{scipy.optimize.differential\_evolution}.
        \item \textbf{Geometria:} Biblioteca \texttt{Shapely} para cálculos exatos de interseção e área de sobreposição.
        \item \textbf{Restrições Suaves (Soft Constraints):}
        Em vez de rejeitar soluções inválidas, aplicamos uma penalidade proporcional ao erro. Isso cria um gradiente que guia o otimizador para regiões válidas.
        \item \textbf{Heurística de Inicialização (4 Cantos):}
        Um algoritmo guloso (Greedy) tenta empacotar as formas partindo dos 4 cantos (BL, BR, TL, TR). A melhor configuração inicializa a população.
    \end{itemize}
\end{frame}

% ============================================================================
% SEÇÃO 4: DISCUSSÃO E RESULTADOS
% ============================================================================
\section{Discussão e Resultados}

\begin{frame}{Tempo de Execução vs. Complexidade (Escalabilidade)}
    \textbf{Método A (Analítico):}
    \begin{itemize}
        \item Extremamente rápido em cenários simples (1 a 6 formas): frações de segundo (0.01s - 0.11s).
        \item \textbf{Salto de complexidade:} Comportamento bimodal. Salta de $\sim 0.1$s ($<10$ formas) para $\sim 50-60$s (15+ formas). Sugere complexidade exponencial.
    \end{itemize}

    \vspace{0.5cm}
    \textbf{Método M (Metaheurística):}
    \begin{itemize}
        \item Consistentemente mais lento. Média de $\sim 41$s mesmo para uma única forma.
        \item Curva de crescimento menos abrupta, mas opera em patamar superior (40s a 200s).
    \end{itemize}
\end{frame}

\begin{frame}{Qualidade da Solução (Área Encontrada)}
    \textbf{Superioridade do Método Analítico (A):}
    \begin{itemize}
        \item Encontrou áreas maiores em \textbf{87.5\%} dos casos testados.
    \end{itemize}

    \vspace{0.3cm}
    \textbf{Degradação do Método M em alta complexidade:}
    \begin{itemize}
        \item \textbf{Poucas formas (1-6):} Resultados próximos ao Analítico (diferença $< 2\%$).
        \item \textbf{Alta complexidade (15-20):} Performance desaba.
        \begin{itemize}
            \item Método A: média $\approx 5500-7000$.
            \item Método M: cai para $\approx 1400-4000$.
        \end{itemize}
    \end{itemize}
    
    \vspace{0.2cm}
    \textit{Indicação:} A Metaheurística tem dificuldade em convergir para ótimos globais em espaços de busca muito restritos.
\end{frame}

\begin{frame}{Conclusão Geral}
    \textbf{Método Analítico (A):}
    \begin{itemize}
        \item Claramente superior neste conjunto de dados.
        \item Melhor eficiência de tempo e qualidade final, especialmente em cenários complexos.
    \end{itemize}

    \vspace{0.5cm}
    \textbf{Metaheurística (M):}
    \begin{itemize}
        \item Viável apenas para instâncias muito simples (empate em qualidade, perde em tempo).
        \item Áreas menores em casos complexos (15+ formas) sugerem prisão em ótimos locais ou tempo insuficiente para exploração.
    \end{itemize}
\end{frame}

\begin{frame}
    \centering
    \Huge \textbf{Obrigado!}
    
    \vspace{1cm}
    \large Perguntas?
\end{frame}

\end{document}
